\subsection{Massive MIMO for 6G}
Multiple-Input Multiple-Output (MIMO) technology is a fundamental building block of 6G systems \cite{bojovic_mimo,cazzella_dt_mimo}.
Large scale MIMO antenna arrays are essential for operating in high-frequency bands such as mmWave and THz \cite{uwaechia_mmwave,han_thz}.
With the evolution of MIMO antenna arrays, they were able to provide highly directional beamforming and increased angular resolution, allowing the network to distinguish users and objects in environment \cite{bojovic_beamforming,gao_isac_mimo}.
This capabilities are important for emerging 6G applications where it has operate in complex, dynamic environment.

In particular, MIMO plays a improtant role in Integrated Sensing and Communication (ISAC), where wireless channel is exploited not only for data transmission but also for environment sensing \cite{zhou_isac,gao_isac_mimo,hua_isac_beamforming}.
In such systems, spatial information of channel such as angles, delays, etc carry information about the physical environment including position and movment of objects.
Capturing accurately of these information is critical for enabling functionalities such as beam tracking, object detection, and context-aware decision-making.

However, testing and deploying such kind of systems in real world is challenging due to high cost, complex setup, and limited availability of testing facilities.
Simulation become important to overcome these limitation for 6G research.
It provide a controlled and scalable environment to evaluate advance MIMO techniques under realistic conditions.
However, conventional stochastic channel models, that are widely use in simulators, are designed for statical performance evaluation and lack geometric information required for application such as ISAC \cite{tr38901,alsamman_mmwave,han_thz}.
This limitation movitates the adoption of physically grounded simulation approaches that can accurately represents the interaction between electromagnetic waves and the environment.
